%% ============================================================
%%  SECCIÓN 4 — Fundamento Teórico
%%  Responsable: Minilux
%% ============================================================

\section{Fundamento Teórico}

\subsection{Fuerzas intermoleculares y tensión superficial}

Las moléculas de un líquido se atraen entre sí mediante fuerzas cohesivas de
corto alcance. Una molécula en el interior del líquido está rodeada
simétricamente por vecinas en todas las direcciones, por lo que la fuerza
resultante neta es nula. En cambio, una molécula ubicada en la superficie
carece de vecinas en la fase gaseosa superior, de modo que la resultante apunta
hacia el interior y es perpendicular a la interfaz. Esta asimetría confiere a
la capa superficial un exceso de energía libre que se manifiesta como una
tendencia a contraerse, dando lugar al fenómeno conocido como \textit{tensión
superficial} \cite{serway}.

\subsection{Coeficiente de tensión superficial}

El \textit{coeficiente de tensión superficial} $\gamma$ se define como la
fuerza por unidad de longitud que ejerce una porción de superficie líquida
sobre la porción adyacente a lo largo de una línea imaginaria trazada sobre
ella \cite{halliday}:
\begin{equation}
    \gamma = \frac{F}{L},
    \label{eq:gamma_def}
\end{equation}
donde $F$ es la componente de fuerza tangente a la superficie y $L$ la
longitud de la línea de contacto. Las unidades en el SI son
\unit{\newton\per\meter}, equivalentes a \unit{\joule\per\meter\squared}, lo
que evidencia que $\gamma$ también representa la energía superficial por unidad
de área. El valor aceptado para el agua pura a \SI{20}{\celsius} es
$\gamma_{\text{ref}} = \SI{72.8e-3}{\newton\per\meter}$ \cite{serway}.

\subsection{Equilibrio de torques — primer método}

La balanza de Mohr--Westphal opera bajo la segunda condición de equilibrio
estático: la suma de torques respecto al eje de rotación es nula,
\begin{equation}
    \sum_{i} \tau_i = 0.
    \label{eq:equilibrio_torques}
\end{equation}
En el primer método, el anillo metálico sumergido en agua ejerce sobre la
superficie una fuerza de tensión superficial $F_{TS}$ a lo largo de su
perímetro de contacto, tanto interno como externo. Esta fuerza actúa en el
brazo mayor de la balanza a una distancia $d_{\text{brazo}}$ del eje. Los
jinetillos de masa $m_j$ colocados en la ranura $n_j$ del brazo mayor (de 10
divisiones totales) generan el torque equivalente que equilibra el sistema.
Aplicando la Ec.~\eqref{eq:equilibrio_torques}:
\begin{equation}
    F_{TS} \cdot d_{\text{brazo}} =
    \left(\sum_{j} m_j \,\frac{n_j}{10}\right) g \cdot d_{\text{brazo}},
    \label{eq:torque_balance}
\end{equation}
de donde se obtiene la masa efectiva equivalente
\begin{equation}
    M_{\text{ef}} = \sum_{j} m_j \,\frac{n_j}{10}
    \label{eq:masa_ef}
\end{equation}
y la fuerza de tensión superficial
\begin{equation}
    F_{TS} = M_{\text{ef}}\, g.
    \label{eq:fuerza_ts}
\end{equation}
La longitud de contacto es la suma de las circunferencias interna y externa del
anillo:
\begin{equation}
    L = 2\pi R_{\text{ext}} + 2\pi R_{\text{int}}
      = 2\pi\,(R_{\text{ext}} + R_{\text{int}}),
    \label{eq:longitud_contacto}
\end{equation}
por lo que, combinando las Ecs.~\eqref{eq:gamma_def},
\eqref{eq:fuerza_ts} y \eqref{eq:longitud_contacto}:
\begin{equation}
    \gamma = \frac{M_{\text{ef}}\,g}{2\pi\,(R_{\text{ext}} + R_{\text{int}})}.
    \label{eq:gamma_m1}
\end{equation}

\subsection{Equilibrio de fuerzas — segundo método}

En el segundo método se forma una película jabonosa entre dos tubitos de vidrio
unidos por un hilo. Al suspender el sistema por el tubo superior, la tensión
superficial de la película curva el hilo inferior, que se modela como un arco
de circunferencia de radio $R$ \cite{halliday}. Las magnitudes geométricas
relevantes son: $a$, el semiancho del tubo; $b$, el semiancho mínimo de la
película curvada; y $h$, la semialtura de la película. 

De la geometría del arco se obtienen las relaciones
\begin{equation}
    a + R\sin\alpha = b + R, \qquad h = R\cos\alpha,
    \label{eq:geometria_arco}
\end{equation}
donde $\alpha$ es el ángulo que forma la tangente al arco en el punto de
contacto con el tubo respecto a la vertical.

Aplicando equilibrio de fuerzas sobre el tubo inferior en la dirección
vertical:
\begin{equation}
    \sum F_y = P - 2\gamma(2a) - 2T\sin\alpha = 0,
    \label{eq:eq_vertical}
\end{equation}
y en la dirección horizontal sobre la película:
\begin{equation}
    \sum F_x = 2\gamma(2h) - 2T\cos\alpha = 0,
    \label{eq:eq_horizontal}
\end{equation}
donde $T$ es la tensión en el hilo y $P = mg$ el peso del tubo inferior.
Dividiendo la Ec.~\eqref{eq:eq_vertical} entre la
Ec.~\eqref{eq:eq_horizontal} y usando las relaciones
geométricas~\eqref{eq:geometria_arco}, se demuestra que:
\begin{equation}
    \gamma = \frac{P}{2\!\left(\dfrac{h^{2}}{a - b} + a + b\right)}.
    \label{eq:gamma_m2}
\end{equation}
La demostración detallada de la Ec.~\eqref{eq:gamma_m2} se desarrolla en el
Anexo~A.

\subsection{Efecto de los surfactantes}

Los surfactantes son moléculas anfifílicas que se adsorben preferentemente en
la interfaz líquido--aire, reduciendo las fuerzas cohesivas entre moléculas de
agua y disminuyendo $\gamma$ de forma apreciable. El detergente empleado en el
segundo método actúa como surfactante, por lo que se espera que el valor de
$\gamma$ obtenido en ese método sea significativamente inferior al del primer
método, realizado con agua pura \cite{serway}.